\section{Related Work}
\label{sec:rw}
% - Machine Unlearning (постановка задачи)
% - Unlearning Benchmarks (кратко пробежаться по TOFU, MUSE, WMDP, подсветив, что они не учитывают частотность/популярность фактов)
% Лламу чаще всего используют для тестирования алгоритмов и бенчмарков в анлернинге

\textbf{Machine Unlearning.}
Machine Unlearning (MU) aims to selectively remove the impact of specific data points from a trained model without requiring a full, costly retraining \cite{mantelero2013eu, cao2015towards, sekhari2021remember}. The goal is to obtain a model that behaves as if the forgotten data was never part of its training set. Initial work focused on traditional machine learning, while recent efforts have shifted to textual unlearning in LLMs \cite{sekhari2021remember, kurmanji2023towards}.
A significant number of modern unlearning approaches and benchmarks are tested primarily on small models from the LLaMA-3 family (up to 8 billion parameters) \cite{yuan2025towards, maini2024tofu, wang2025selective, si2023knowledge}, as they have shown high effectiveness in unlearning tasks \cite{maini2024tofu, yuan2025towards}. For this reason, our experiments are also centered around this model architecture.
The unlearning algorithms used in this research are described in more detail in Section \ref{sec:train}.

\textbf{Unlearning Benchmarks.}
There are several benchmarks available for evaluating MU algorithms. For instance, some benchmarks rely on synthetic data to create controlled but artificial scenarios \cite{geng2025comprehensive, hu2024unlearning, si2023knowledge}. A prominent example is TOFU \cite{maini2024tofu}, which consists of 200 fictitious author profiles to simulate the unlearning of specific, non-public facts. TOFU is one of the most popular modern MU benchmarks, they use 1\%, 5\% and 10\% of the full benchmark for forgetting, the rest is used to evaluate the quality of unlearning. While useful for testing privacy-like scenarios, this synthetic nature does not capture the complexity of unlearning real-world knowledge of varying prevalence.

Other benchmarks use real-world data but concentrate on evaluation objectives orthogonal to fact popularity. For instance, WMDP \cite{li2024wmdp} includes 3,668 multiple-choice questions specifically designed to evaluate the unlearning of hazardous knowledge in LLMs. Similarly, MUSE \cite{shimuse} offers a multifaceted evaluation of unlearning, assessing properties like efficacy and efficiency. Although vital, the primary goals of these benchmarks are not to investigate how the inherent properties of knowledge, like its familiarity, affect unlearning.

